% !PRINTED 372
\weberpage{372}
If we form the two forms
\[
\begin{aligned}
\lambda &= a_1 a_2\,t_1 + \frac{b-\sqrt D}{2}\,t_2, \\
\lambda' &= a\,t_1 + \frac{b-\sqrt D}{2}\,t_2,
\end{aligned}
\]
then it follows that:
\[
\begin{aligned}
N(\lambda) &= a_1 a_2 (a_1 a_2\,t_1^2 + b\,t_1 t_2 + c\,t_2^2), \\
N(\lambda') &= a (a\,t_1^2 + b\,t_1 t_2 + a'c\,t_2^2).
\end{aligned}
\]
The two forms $(a_1 a_2,\ b,\ c)$, $(a,\ b,\ a'c)$ are primitive (because a common prime divisor of $a_1 a_2$, $b$, $c$ would, on account of (15), have to divide $Q$), and for the absolute norm it follows that:
\[
\begin{aligned}
N_a(\lambda) &= a_1 a_2 = N(\mathfrak{a}\mathfrak{b}), \\
N_a(\lambda') &= a;
\end{aligned}
\]
but since $\lambda'$ is divisible by $\mathfrak{a}\mathfrak{b}$, $a$ must be divisible by $a_1 a_2$, and consequently $a'=1$.

From this it follows:

3. \gesperrt{If $a_1$, $a_2$, $b$ are without common divisor and}
\[
D = b^2 - 4\,a_1 a_2 c,
\]
\gesperrt{then the form}
\[
T = (a_1 a_2,\ b,\ c)
\]
\gesperrt{is compounded from the two forms:}
\[
U = (a_1,\ b,\ c\,a_2), \quad V = (a_2,\ b,\ c\,a_1)
\]

Applying this to two opposite forms
\[
U = (a,\ b,\ c), \quad V = (c,\ b,\ a)
\]
there results
\[
T = (ac,\ b,\ 1),
\]
and this is equivalent to the principal form $\left(1,\ 0,\ \dfrac{-D}{4}\right)$ or $\left(1,\ 1,\ \dfrac{1-D}{4}\right)$.

4. \gesperrt{Two opposite classes, compounded, give the principal class.}

In order to compose a class $A$ repeatedly with itself, one chooses $\mathfrak{a}$ in $A$ so that it is relatively prime to $D$, and denotes by $(a,\ b,\ c)$ the corresponding form. Then $a$ is relatively prime to $b$. For if $d$ were a divisor of $a$ and $b$, then $d$ would also be contained in $D$, hence relatively prime to $\mathfrak{a}$. Then $a:d$ too would be divisible by $a$, which, if $d>1$, contradicts the definition of $a$.
